% Chapter 6 -- Reproducible builds.  Sources (Law 1): ADR-0007 DEC-04
% (insertion-ordered tables, no hashing; one label counter, never
% re-derived; no timestamps, host paths or nondeterministic iteration),
% DEC-06 (reproducibility is the precondition; test-seed-repro);
% os/tps/CODEGEN.md s6; CLAUDE.md Rule 11.  Proved by test-seed-repro
% (two scratch builds compared) and test-tps-gen-fpc (every corpus program
% and tps.pas generated twice, compared byte for byte).
\Gates{test-seed-repro, test-seed-repro-mutant, test-tps-gen-fpc, test-tps-gen-mutant, test-pubs-turbo}
\chapter{Reproducible builds}
\ix{reproducibility}

Every step of the bootstrap must give the same bytes from the same input,
on any development system, at any time.  Without that, the comparison of
$K_2$ and $K_3$ proves nothing.

\section{The rules}
Turbo Initech, the genesis compiler and the build that runs them keep these
rules:

\begin{itemize}
\item \textbf{No date, time or place in any output.}  No compiler, assembly
  file or program records when, where or on which machine it was made.
\item \textbf{Tables in declaration order.}  Every table of the compiler is
  an array kept in the order in which the declarations appear.  No table is
  a hash table, whose order could depend on anything but the program.
\item \textbf{One counter for labels.}  Every label takes its number from a
  single counter that runs through the whole compilation; no second pass
  computes label numbers again.
\item \textbf{One output for one input.}  The compiler has no options and
  no settings: its output depends only on the bytes of \file{TPSIN.PAS}.
\end{itemize}

\section{The checks}
\ix{reproducibility, checks}
The build checks the rules directly.  Procedure TIN-P-01 builds every test
program of the genesis compiler twice, in two separate directories, and
compares the assembly language, the objects and the linked programs.
Procedure TIN-P-03 compiles every corpus program and the Turbo Initech
source twice and compares the results, and its first lines report each
program as \code{deterministic}.  Each of these checks is itself checked:
a deliberately broken genesis compiler that varies its output, and a
deliberately broken Turbo Initech that reuses label numbers, must each be
caught by the build.

\section{When a build is not reproducible}
A difference between two builds of the same source is a defect, never a
nuisance to be worked around.  Look first for a value that depends on the
time, the directory, the order of files in a directory, or memory that
was never set; then for a table read in an order that is not the order of
the source.  Report the two differing files with the defect.
